
\section{Introduction}\label{sec:introduction}
Fix \(0<\rho<1\), the chamber weights \((6,6\sqrt2,6\sqrt2,6)\),
and an odd integer \(L\ge4^8\), an integral power \(M\) of \(L\) with
\(M\ge L^2\), integers \(r\ge2\) and \(p_0,p>0\), \(0<\varepsilon\le1/200\),
\(1-\varepsilon<\alpha<1\), \(c_{\rm fl}\ge1\), and \(0<\bar m\le1\).
The photon is periodic and the four-component chamber spinors are
antiperiodic in all three directions. Each photon site has 48 chamber
labels, with measure \(\eta_k^3/48\) at mesh \(\eta_k=L^{-k}\).
For an integer \(N\ge0\), set
\[
 h=L^{\lfloor\log_L(1/\bar m)\rfloor},\qquad
 \eta=L^{-N},\quad K=N+\log_Lh,\quad
 a_j=L^{j-N},\quad e_j=e\sqrt{a_j},\quad \mu_j=\bar m a_j.
\]
Here \(0\le j\le K\), and \(q=e\sqrt\eta\) denotes the charge at
level zero. We write \(\bar e_0,\bar e_1\) for the fixed charge
ceilings, to distinguish them from the running charges \(e_j\).
Put \(v_j=-\log e_j\), and let \(r_j\) be the least power \(L^i\),
\(i\in\mathbb Z_{\ge0}\), with \(r_j\ge v_j^r\). The physical sides are
\[
 R_{\rm top}=hMr_K,\quad R^*=R_{\rm top}/L,\qquad
 \mathcal T(e)=\{hL^T:T\in\mathbb N\cup\{0\},\ hL^T\le R_{\rm top}\}.
\]
They are independent of \(N\). For \(R\in\mathcal T(e)\), define
\(\ell_j=R/a_j\) and
\(\mathfrak n_j=\log_L(\ell_j/(Mr_j))\). The stopping level \(k_*(R)\)
is the first \(j\) with \(\mathfrak n_j=0\), when it exists.
Write \(k_{\rm top}=k_*(R_{\rm top})\). A star denotes evaluation
at \(k_*(R)\): for example \(a_*=a_{k_*(R)}\) and
\(\mu_*=\mu_{k_*(R)}\).

\paragraph{Numerical hypotheses.}
Assume
\[
 \begin{gathered}
 0<e\sqrt h\le\bar e_1\le\min(\bar e_0,1/3),\qquad
 \bar e_0^\varepsilon\le1/4,\\
 c_{\rm fl}L^{-3/8+6/200}\le1/16,\qquad L^{-5/8}\le1/4,
 \end{gathered}
\]
and, for every \(v\ge-\log\bar e_1\),
\begin{equation}\label{eq:cv-scale-conditions}
 \begin{gathered}
 v\ge8,\qquad
 \left(1+\frac{3\log L}{2v}\right)^r<L,\qquad
 v^r>31L+16,\\
 (c_{\rm fl}+2)e^{-\varepsilon v}
       \le\frac{1-L^{-\varepsilon/2}}{16}.
 \end{gathered}
\end{equation}

\paragraph{The activity maps.}
Use the real chamber activity spaces, with independent connection
arguments, defined in \cref{sec:activities}. Write \(\mathscr E_j\) for an activity
with external connection \(\mathscr C\), and \(E_j=\mathscr E_j|_{\mathscr C=0}\).
At every used scale impose all the simultaneous hypotheses
(i)--(iv) of \cite[Theorem~A]{NonlinearStep}, including its
primitive, inverse-domain, determinant, extraction, reconstruction,
contour and connected-logarithm conditions. The constants, norms and
operators are those of that theorem and of
\cite[Theorems~A--B]{NormalizedContraction}; their definitions and the
parts used below are recalled in \cref{sec:activities,sec:inputs}.
These are hypotheses on the actual chamber maps. In particular the
following tuning inequalities do not replace their domain or numerical
conditions. The decay weight satisfies \(\kappa\ge7\kappa_0+7\), where
\(\kappa_0=54\log72\); the constants are calibrated together as in the
cited step theorem.

On the top torus choose the photon and fermion families of
\cref{subsec:rgcov-top-system}, with polymer decay
\(\kappa_w+\kappa_0\), where \(\kappa_w=\kappa+\kappa_0+3\).
On each smaller odd quotient use their projections, with common
constant \(C_w=64K_0C_w^{(1)}\), \(K_0=\exp(\kappa_0)/36\).
Both incoming and outgoing sides are required to have
\(\ell/(LM)\) odd and at least 13. The step comparison requires
\(\ell_{\rm small}/L>(31L+16)M\); the stopping calculation verifies
these side conditions at all steps used in the theorem.

For a comparison with smaller unit side $\ell_1$, a connected
covering polymer $\widetilde X$ is admissible when its smallest
rectangular box has largest side at most $\ell_1-14(2L+1)M$.
Its two torus images are called flat. Local agreement means equality
under the common covering generators for connection pairs agreeing
on $\operatorname{box}(\widetilde X)^{\sim6(2L+1)}$.
Here the superscript denotes enlargement by that number of $M$-cube
layers. \Cref{prop:rgcov-common-lifts} justifies this identification.

The estimates used for tuning take the following form.
Set \(w_j^{\mathrm m}=e_j^{3/4-8\varepsilon}\) and
\(w_j^{\mathrm E}=e_j^{1/4-7\varepsilon}\).
The no-slot data norm is
\(\max\{|m|/w_j^{\mathrm m},\|E\|_j/w_j^{\mathrm E}\}\); with an external slot,
replace \(\|E\|_j\) by \(\|\mathscr E\|_j^x\).
Radius-four data balls mean the corresponding closed balls in the
complexified spaces. Unit balls below are closed balls in their real
subspaces. At zero the activity coordinate is fixed to zero.
The maps preserve reality and obey
\begin{align}
 m_{j+1}&=L(m_j+m(E_j)),&
 \varepsilon_{j+1}^v&=L^3(\varepsilon_j^v+\varepsilon_v(E_j)+\delta\varepsilon_j^0),
 \label{eq:cv-scalar-recursion}\\
 E_{j+1}&=\Phi_j(m_j,E_j),&
 \Phi_j(m,E)&=\mathcal L(RE)^{\rm nat}
                  +\mathcal LE_j^{\det}+E_j^*(m,E).
 \label{eq:cv-activity-recursion}
\end{align}
Restriction to \(\mathscr C=0\) commutes with the full external-slot recursion
and does not increase its norm. Scalar extraction is linear and real,
\[
 |m(E)|\le c_{\rm fl}e_j^{1/2}\|E\|_j,\qquad
 |\varepsilon_v(E)|\le c_{\rm fl}\|E\|_j,\qquad
 0\le\delta\varepsilon_j^0\le e_j^7.
\]
On the unit ball the normalized size and difference estimates are
\begin{equation}\label{eq:cv-normalized-inputs}
 \begin{aligned}
 L^{-1}w_{j+1}^m/w_j^{\mathrm m}&\le1/4,\\
 \|\Phi_j(m,E)\|_{j+1}/w_{j+1}^E
   &\le\tfrac1{16}\|E\|_j/w_j^{\mathrm E}+2e_{j+1}^{\varepsilon},\\
 |m(E)-m(\widetilde E)|/w_j^{\mathrm m}
   &\le c_{\rm fl}e_j^\varepsilon d_E,\\
 \|\Phi_j(m,E)-\Phi_j(\widetilde m,\widetilde E)\|_{j+1}/w_{j+1}^E
   &\le\tfrac1{16}d_E+e_{j+1}^\varepsilon(d_m+d_E),
 \end{aligned}
\end{equation}
where \(d_m=|m-\widetilde m|/w_j^{\mathrm m}\) and
\(d_E=\|E-\widetilde E\|_j/w_j^{\mathrm E}\).
The same size estimate holds for the external-slot norm. The
fixed-point argument uses the no-slot difference estimate.
Initially \(E_0=\mathscr E_0=\mathfrak P_0=0\), where
\(\mathfrak P_j\) is the scalar external family in \cref{sec:inputs}.
The first step uses the chamber definitions in \cref{sec:inputs}.

\begin{theorem}[A common mass and vacuum pair]\label{thm:main}
Under the hypotheses above, the following assertions hold.
For each \(R\in\mathcal T(e)\) there is a finite threshold
\(N_{\min}(R,e)\) such that the first stop exists for
\(N\ge N_{\min}(R,e)\). Its physical quantities are then independent
of \(N\). The thresholds are nonincreasing and the stopping levels
nondecreasing
with \(R\); their precise formulas are in \eqref{eq:rgfl-mesh-threshold}.

For \(N\ge N_{\min}(R_{\rm top},e)\), there is exactly one real
pair \(P_N(e)=(m_{0,N}(e),\varepsilon_{0,N}^v(e))\) whose top-torus
trajectory is controlled and has
\(m_{k_{\rm top}}=\varepsilon_{k_{\rm top}}^v=0\).
Controlled means that each preceding step is defined in its
radius-four data ball and \(|m_j|\le w_j^{\mathrm m}\) through the stop.
For every \(R\in\mathcal T(e)\) with \(N\ge N_{\min}(R,e)\),
the same pair is controlled on \(R\). Up to its stop, the scalar
mass and vacuum sequences equal the top sequences, while the
activities and external scalar families agree on admissible flat
lifts under the algebra identification of \cref{prop:rgcov-common-lifts}. On each such torus,
\begin{equation}\label{eq:cv-common-bounds}
 |m_j|\le\tfrac18e_j^{3/4-8\varepsilon},\qquad
 \|E_j\|_j\le\|\mathscr E_j\|_j^x\le4e_j^{1/4-6\varepsilon},\qquad
 |\varepsilon_j^v|\le5c_{\rm fl}e_j^{1/4-6\varepsilon}.
\end{equation}
The smaller terminal couplings need not vanish.

Within the fixed construction, the pair is determined by
\((N,e,\bar m)\) and is independent of the smaller side, twist,
external connection and probe. Define only \(m_{0,N}(e)=0\) below
the top definition threshold. For every \(N\ge0\),
\[
 m_{0,N}(e)\in\mathbb R,\qquad
 |m_{0,N}(e)|\le\tfrac18q^{3/4-8\varepsilon}<\tfrac18 .
\]
The total bare mass \(\boldsymbol\mu_N=\bar mL^{-N}+m_{0,N}(e)\)
has no asserted sign. On \(R^*\), for
\(N\ge N_{\rm ref}(e):=N_{\min}(R^*,e)\),
\[
 a_*\in\{h/L,h/L^2\},\qquad
 \ell_*=Mr_*,\qquad r_*>31L+16,\qquad
 L^{-3}<\mu_*\le L^{-1}.
\]


\end{theorem}

The theorem concerns the prescribed family of odd reference tori.
The scalar pair is selected before inserting a smaller side, a harmonic
twist or an external probe. The proof has two parts. On the top torus,
the mass equation is solved backwards from zero terminal mass, while
the activity equation runs forwards from zero initial activity. In the
weighted sequence norm these two equations define a contraction with
factor \(5/16\). The vacuum is recovered by an explicit backward sum.

To use that pair on a smaller torus, we compare the actual local
operations. Odd covers preserve the antiperiodic character. Projection
gives a signed finite image sum; when weakening parameters confine a
row to a common nonwrapping lift, only one image survives. The local
Grassmann algebra then identifies extraction, determinant traces, both
finite localizations and the fluctuation integrals. Absolute convergence
passes this equality to the connected logarithm. At each smaller-volume
step, its own continuation estimate controls the full activity norm,
and local comparison identifies the scalar couplings with those of the
top trajectory.

The cover and tuning arguments are those of
\cite[\S\S15--16]{QED3Master}, based on
\cite[Chapters~27--28]{ChamberRG}. The stopping-scale calculation is
the one retained there from \cite[Chapter~19]{LatticeCurvature}.
The photon, fermion, extraction and nonlinear constructions are supplied
by \cite{MaxwellBlockManuscript,FermionBlockManuscript,NormalizedContraction,NonlinearStep}.
We give the cover comparison and finite-sequence tuning proofs here,
with those constructions as explicit inputs.

The definitions in \cref{sec:activities} and the input identities in
\cref{sec:inputs} fix the maps compared in
\cref{sec:rg-cover-comparison,sec:local-comparison}.
The stopping calculation is in \cref{sec:rg-flow};
\cref{sec:common-tuning} proves the common-pair theorem.
